\documentclass[12pt]{article}
\usepackage[utf8]{inputenc}
\usepackage[T1]{fontenc}
\usepackage{hyperref}
\title{Yendang Language Resources for K-12 Curriculum Development}
\author{Compiled by Claude Sonnet 5.5 (claude-sonnet-5-5) for Andr\'as Kornai}
\date{}
\begin{document}
\maketitle

\section{Introduction}
Yendang (ISO 639-3: ynq; Glottocode yend1241) is spoken in northeastern Nigeria (Adamawa and Taraba states). Wikipedia places it in the Leko--Nimbari group of Savanna languages. Its code changed from \emph{yen} to \emph{ynq} in 2012. Listed dialects are Kuseki, Yofo and Poli (Akule, Yakule); Joshua Project also gives Akule and Ikume as names. Speaker estimates differ: about 50,000 (1987, Wikipedia and ASJP) against 133,000 for the Yandang people group (Joshua Project, from Ethnologue 2016 data). The searches found no school materials in the language. What exists is reference-level information and missionary audio material.

\section{Primary Educational Resources}
No primers, textbooks or school programmes for Yendang turned up in the searches. This section is empty; a curriculum would have to be built from scratch with the community.

\section{Linguistic and Reference Materials}
\begin{itemize}
\item Wikipedia, ``Yendang language'': overview of classification, dialects and speaker numbers. \url{https://en.wikipedia.org/wiki/Yendang_language}
\item ASJP database entry for Yendang, a wordlist and speaker figure. \url{https://asjp.clld.org/languages/YENDANG}
\end{itemize}
No grammar or dictionary was found in the search results.

\section{Religious and Community Resources}
\begin{itemize}
\item Joshua Project, Yandang in Nigeria: people-group profile. It reports Bible portions translated in 2016, with no New Testament or full Bible. \url{https://joshuaproject.net/people_groups/15978}
\item Joshua Project language page for ynq. \url{https://joshuaproject.net/index.php/languages/ynq}
\end{itemize}

\section{Digital and Audio Resources}
\begin{itemize}
\item Global Recordings Network, ``Yendang: Akule \&amp; Ikume'': audio Bible-story recordings. \url{https://globalrecordings.net/en/language/1878}
\end{itemize}

\section{Curriculum Development Context}
The language appears to be vigorous but has no known literacy materials or school use. Any curriculum would need an agreed orthography, community and language-committee involvement, and new materials. The speaker figures should be verified because sources disagree widely. Nothing here is a ready-made classroom resource.

\end{document}